% !TEX root =./ECE444_Practice_main.tex
%
% L9 practice set (Loop and Monopole Antennas) -- image theory, the quarter-wave
% monopole over real ground, the electrically small loop, and the size-bandwidth trade.

\fullwidth{\textbf{Learning Objective 2.1: } I can describe the radiation behavior of
simple resonant antennas (isotropic radiator, half-wave dipole, monopole, loop) and
calculate their gain and impedance.
\\[4pt]\normalfont\small Show your work. A \textbf{Documentation} line at the end
is required (write \textbf{None} if you did not collaborate).}

\begin{questions}

\question \textbf{A whip for the guard frequency.} You are asked to build a quarter-wave
monopole for $f = 121.5\ \mhz$ on a large aluminum ground plane, fed with $50\ \ohms$ coax.
Take the thin-wire half-wave dipole as $73 + j42.5\ \ohms$.
\begin{parts}

	\begin{minipage}[t]{0.58\linewidth}
		\part How long is the whip?
		\begin{solution}[1.0in]
		\eq{ \lambda &= \frac{c}{f} = \frac{3\ee{8}}{121.5\ee{6}} = 2.47\ \m \\
		     L &= \frac{\lambda}{4} = 0.617\ \m }
		About 62 cm of metal above the plane.
		\end{solution}
	\end{minipage}\hfill%
	\begin{minipage}[t]{0.37\linewidth}
		\raggedleft \vspace{12pt}
		$L$ = \ansbox[1.3in]{$0.617\ \m$}
	\end{minipage}

	\part State the input impedance you expect at the base of the whip, and explain in one
	or two sentences why it is not the dipole value.
		\begin{solution}[1.1in]
		$Z_{\text{in}} = 36.5 + j21.3\ \ohms$, exactly half the dipole value. The image in the ground
		plane supplies the missing lower half of the dipole, so the current at the feed is the
		same as the dipole's, but the generator only drives half the structure and therefore
		only has to supply half the voltage. Half the voltage at the same current is half the
		impedance.
		\end{solution}

	\begin{minipage}[t]{0.58\linewidth}
		\part Find $\vert \Gamma \vert$ and the VSWR seen by the $50\ \ohms$ feed.
		\begin{solution}[1.3in]
		\eq{ \Gamma &= \frac{Z_{\text{in}} - Z_0}{Z_{\text{in}} + Z_0}
		            = \frac{-13.5 + j21.3}{86.5 + j21.3} \\
		     \vert \Gamma \vert &= \frac{25.2}{89.1} = 0.283 \\
		     \text{VSWR} &= \frac{1 + 0.283}{1 - 0.283} = 1.79 }
		Return loss $= -20\mylog{0.283} = 11.0\db$.
		\end{solution}
	\end{minipage}\hfill%
	\begin{minipage}[t]{0.37\linewidth}
		\raggedleft \vspace{12pt}
		VSWR = \ansbox[1.3in]{$1.79$}
	\end{minipage}

	\part You may change the whip or the ground plane, but you may not add parts. Give two
	changes that improve the match, and say what each one does to the impedance.
		\begin{solution}[1.4in]
		(1) \emph{Trim the whip} to about $0.24\lambda$ (roughly 59 cm). Shortening it a few
		percent tunes out the $+j21.3\ \ohms$ of reactance and leaves about $36\ \ohms$ real,
		so VSWR drops to $50/36 = 1.39$. The reactance, not the resistance, is what costs most
		of the mismatch.
		(2) \emph{Droop the radials} of the ground plane down about $45\mydeg$. Tilting the
		radials raises the base resistance from roughly $36\ \ohms$ toward $50\ \ohms$, giving
		VSWR near 1.0 with no added components. This is why commercial ground-plane antennas
		have sagging radials.
		\end{solution}

\end{parts}

\newpage

\question \textbf{Which way do you lay the wire?} A perfectly conducting deck acts as an
image plane. An antenna at height $h$ above it radiates as the element pattern times a
two-element array factor, $2\vert\cosp{kh\cos\theta}\vert$ for a vertical element and
$2\vert\sinp{kh\cos\theta}\vert$ for a horizontal one, with $\theta$ measured from the
vertical.
\begin{parts}

	\begin{minipage}[t]{0.58\linewidth}
		\part A horizontal wire sits at $h = 0.05\lambda$. By how many dB is its
		straight-up ($\theta = 0$) field below that of the same wire at $h = 0.25\lambda$?
		\begin{solution}[1.3in]
		\eq{ h = 0.05\lambda: \quad 2\sinp{2\pi(0.05)} &= 2\sinp{18\mydeg} = 0.618 \\
		     h = 0.25\lambda: \quad 2\sinp{2\pi(0.25)} &= 2\sinp{90\mydeg} = 2.00 \\
		     20\mylog{0.618/2.00} &= -10.2\db }
		\end{solution}
	\end{minipage}\hfill%
	\begin{minipage}[t]{0.37\linewidth}
		\raggedleft \vspace{12pt}
		$\Delta$ = \ansbox[1.3in]{$-10.2\db$}
	\end{minipage}

	\begin{minipage}[t]{0.58\linewidth}
		\part Evaluate the vertical array factor at the horizon ($\theta = 90\mydeg$) for
		an arbitrary height $h$. What does the result tell you?
		\begin{solution}[1.1in]
		\eq{ 2\vert\cosp{kh\cos 90\mydeg}\vert = 2\vert\cosp{0}\vert = 2 \quad
		     \text{for every } h }
		A vertical element and its image always add in phase along the horizon, at any height,
		including zero. That is why a vertical antenna works while standing on the ground.
		\end{solution}
	\end{minipage}\hfill%
	\begin{minipage}[t]{0.37\linewidth}
		\raggedleft \vspace{12pt}
		AF = \ansbox[1.3in]{$2$}
	\end{minipage}

	\part In one or two sentences, explain why a horizontal wire lying directly on a metal
	deck radiates essentially nothing.
		\begin{solution}[1.0in]
		A horizontal (tangential) current images with the opposite sign, so the deck creates
		an equal and opposite current directly beneath the wire. As the height goes to zero
		the two are collocated and cancel everywhere, and the radiation resistance collapses
		with them. The wire becomes a heater, not an antenna.
		\end{solution}

	\part You must put up a field-expedient antenna on a large metal roof for
	horizon-to-horizon communication, and you can only get about $0.05\lambda$ of height.
	Which orientation do you choose, and why?
		\begin{solution}[1.2in]
		Vertical. At $0.05\lambda$ a horizontal wire is fighting its own image and loses more
		than 10 dB of the little it radiates, and perfect ground puts a null on the horizon for
		horizontal polarization at \emph{every} height, which is exactly where the traffic is.
		A vertical element gets the in-phase image for free, keeps its pattern maximum on the
		horizon, and needs no height at all.
		\end{solution}

\end{parts}

\newpage

\question \textbf{A magnetic loop, costed honestly.} A single-turn circular loop of radius
$a = 0.30\ \m$ is made of copper wire of radius $b = 3\ \text{mm}$ and operated at
$f = 14.2\ \mhz$. At that frequency the copper surface resistance is
$R_s = 0.98\ \text{m}\ohms$ per square. Recall $R_r = 20\pi^2 \lp C/\lambda \rp^4$ and
$R_{ohmic} = \lp C/2\pi b \rp R_s$ for a single turn.
\begin{parts}

	\begin{minipage}[t]{0.58\linewidth}
		\part Compute $C/\lambda$ and confirm the loop is electrically small.
		\begin{solution}[1.1in]
		\eq{ \lambda &= \frac{3\ee{8}}{14.2\ee{6}} = 21.1\ \m \\
		     C &= 2\pi a = 1.885\ \m \\
		     C/\lambda &= 0.0892 }
		Below the $0.1\lambda$ rule of thumb, so the current is essentially uniform around
		the loop and it behaves as a magnetic dipole.
		\end{solution}
	\end{minipage}\hfill%
	\begin{minipage}[t]{0.37\linewidth}
		\raggedleft \vspace{12pt}
		$C/\lambda$ = \ansbox[1.3in]{$0.0892$}
	\end{minipage}

	\begin{minipage}[t]{0.58\linewidth}
		\part Find the radiation resistance.
		\begin{solution}[1.1in]
		\eq{ R_r &= 20\pi^2 (0.0892)^4 = 197.4 \lp 6.34\ee{-5} \rp \\
		         &= 0.0125\ \ohms = 12.5\ \text{m}\ohms }
		\end{solution}
	\end{minipage}\hfill%
	\begin{minipage}[t]{0.37\linewidth}
		\raggedleft \vspace{12pt}
		$R_r$ = \ansbox[1.3in]{$12.5\ \text{m}\ohms$}
	\end{minipage}

	\begin{minipage}[t]{0.58\linewidth}
		\part Find the ohmic resistance and the radiation efficiency. Express the efficiency
		in dB and give the gain in dBi.
		\begin{solution}[1.5in]
		\eq{ R_{ohmic} &= \frac{C}{2\pi b} R_s = \frac{a}{b} R_s \\
		               &= 100 \lp 0.98\ \text{m}\ohms \rp = 0.098\ \ohms \\
		     \eta_{\text{rad}} &= \frac{R_r}{R_r + R_{ohmic}} = \frac{0.0125}{0.1105} = 0.113 \\
		     10\mylog{0.113} &= -9.5\db \\
		     G &= 1.76 - 9.5 = -7.7\dbi }
		The loop's own copper burns nearly eight times what it radiates.
		\end{solution}
	\end{minipage}\hfill%
	\begin{minipage}[t]{0.37\linewidth}
		\raggedleft \vspace{12pt}
		$\eta_{\text{rad}}$ = \ansbox[1.3in]{$11.3\%$}
		\\[16pt]
		$G$ = \ansbox[1.3in]{$-7.7\dbi$}
	\end{minipage}

	\part The same loop is a respectable receiving antenna at 14.2 MHz but a poor
	transmitting one. Explain the difference in one or two sentences.
		\begin{solution}[1.1in]
		On transmit, 89\% of your power heats the wire, and the tiny radiation resistance
		forces enormous circulating current. On receive, what matters is signal-to-noise, and
		below VHF the external atmospheric noise picked up by the antenna dominates the
		receiver's own noise -- attenuating signal and sky noise together by 9.5 dB leaves the
		ratio essentially unchanged.
		\end{solution}

\end{parts}

\newpage

\question \textbf{Loop or dipole?} Continue with the loop of the previous problem.
\begin{parts}

	\part A direction-finding team rotates a small loop to take a bearing. Why do they turn
	the loop for a \emph{null} rather than for a peak, and where does that null point
	relative to the loop?
		\begin{solution}[1.2in]
		The small-loop pattern is $\vert F \vert = \sin\theta$ with the null along the loop
		\emph{axis} -- straight through the hole -- and the maximum in the plane of the loop.
		A null is sharp: the response changes quickly with a degree of rotation, so the
		bearing is precise. The peak is broad and flat, so a peak-seeking bearing would be
		several degrees wide. When the null is deepest, the transmitter lies along the loop
		axis.
		\end{solution}

	\part A vertical half-wave dipole and a small horizontal loop are both mounted on a mast.
	What polarization does each radiate toward a distant receiver on the horizon?
		\begin{solution}[1.1in]
		The vertical dipole radiates $E_\theta$ -- vertical polarization. The horizontal loop
		is a magnetic dipole pointing straight up; its field is $E_\phi$, which lies parallel
		to the ground at the horizon, so it radiates horizontal polarization. Same donut
		pattern, orthogonal polarizations: cross them and the polarization loss factor from
		Lesson 3 takes the link apart.
		\end{solution}

	\begin{minipage}[t]{0.58\linewidth}
		\part Rewind the loop with $N = 20$ turns. Radiation resistance scales as $N^2$ and
		ohmic resistance as $N$. Find the new efficiency.
		\begin{solution}[1.3in]
		\eq{ R_r &= 20^2 \lp 0.0125 \rp = 5.00\ \ohms \\
		     R_{ohmic} &= 20 \lp 0.098 \rp = 1.96\ \ohms \\
		     \eta_{\text{rad}} &= \frac{5.00}{6.96} = 0.718 }
		71.8\%, a loss of only $1.4\db$: gain $= 1.76 - 1.4 = 0.4\dbi$.
		\end{solution}
	\end{minipage}\hfill%
	\begin{minipage}[t]{0.37\linewidth}
		\raggedleft \vspace{12pt}
		$\eta_{\text{rad}}$ = \ansbox[1.3in]{$71.8\%$}
	\end{minipage}

	\part If turns are that effective, why is a 20-turn loop still not the antenna you would
	choose to transmit with? Name two reasons.
		\begin{solution}[1.3in]
		(1) The $N^2$ scaling assumes the turns do not interact. Packed turns raise the loss
		through proximity effect and distributed capacitance, so the real efficiency falls
		short of the ideal number. (2) Adding turns adds inductance and stored energy without
		making the antenna electrically bigger, so the Q climbs and the usable bandwidth
		shrinks -- and the tuning capacitor sees kilovolts. A resonant loop with
		$C \approx 1\lambda$ solves all of it: about $100$ to $130\ \ohms$ of real,
		radiating impedance and roughly 3.1 dBi, with its maximum along the axis instead of
		in the plane.
		\end{solution}

\end{parts}

\newpage

\question \textbf{The price of being small.} Stay with the single-turn loop of radius
$a = 0.30\ \m$ at $14.2\ \mhz$ ($\eta_{\text{rad}} = 0.113$). The Chu limit from Lesson 3 gives the
radiation quality factor of any antenna that fits inside a sphere of radius $a$ as
$Q \gtrsim 1/(ka)^3 + 1/(ka)$, and the fractional bandwidth is roughly $1/Q$.
\begin{parts}

	\begin{minipage}[t]{0.58\linewidth}
		\part Compute $ka$, the Chu $Q$, and the resulting bandwidth in kHz.
		\begin{solution}[1.4in]
		\eq{ k &= \frac{2\pi}{21.1} = 0.297\ \text{rad/m}, \quad ka = 0.0892 \\
		     Q &\gtrsim \frac{1}{(0.0892)^3} + \frac{1}{0.0892} = 1409 + 11 = 1420 \\
		     \text{FBW} &\approx \frac{1}{Q} = 7.0\ee{-4} = 0.070\% \\
		     \Delta f &= 0.00070 \lp 14.2\ \mhz \rp = 10\ \khz }
		\end{solution}
	\end{minipage}\hfill%
	\begin{minipage}[t]{0.37\linewidth}
		\raggedleft \vspace{12pt}
		$Q$ = \ansbox[1.2in]{$\approx 1420$}
		\\[16pt]
		$\Delta f$ = \ansbox[1.2in]{$10\ \khz$}
	\end{minipage}

	\part A single-sideband voice channel needs about 3 kHz. Does the loop pass it, and what
	does your answer imply about operating across a 350 kHz amateur band?
		\begin{solution}[1.2in]
		Yes, 10 kHz comfortably covers one 3 kHz channel -- the loop is not too narrow to use,
		it is too narrow to leave alone. Moving anywhere else in a 350 kHz band means retuning
		the loop capacitor, because the match walks off within a channel or two. Small antennas
		buy their size back in operator workload.
		\end{solution}

	\part Measured on a VNA, this loop shows a $2{:}1$ VSWR bandwidth several times wider
	than your Chu estimate. Why? Is that good news?
		\begin{solution}[1.3in]
		Loss lowers Q. The Chu number is the \emph{radiation} Q; the loaded Q is reduced
		roughly in proportion to the efficiency, $Q_{loaded} \approx \eta_{\text{rad}} Q \approx
		0.113(1420) \approx 160$, which widens the match to about $0.6\%$, or 90 kHz. It is
		not good news: the extra bandwidth is bought by dissipating your power in copper. A
		$50\ \ohms$ resistor is perfectly matched over every band on the dial and radiates
		nothing.
		\end{solution}

	\part State the size-bandwidth trade as a rule of thumb, in one sentence.
		\begin{solution}[1.0in]
		Shrinking an antenna below about $\lambda/2\pi$ in radius drives the stored near-field
		energy far above the radiated energy, so Q rises roughly as $1/(ka)^3$ and the usable
		bandwidth falls with the cube of the size -- while the radiation resistance falls too,
		so you lose efficiency and bandwidth together.
		\end{solution}

\end{parts}

\vspace{1.5em}
\noindent\textbf{Documentation:} \rule[-2pt]{0.6\linewidth}{0.4pt}

\end{questions}
